\documentclass[a4paper]{article}

% Check for the existence of the variable \pdfoutput:
\newif\ifpdf
\ifx\pdfoutput\undefined
\pdffalse % we are not running PDFLaTeX
\else
\pdfoutput=1 % we are running PDFLaTeX
\pdftrue
\fi

\usepackage{mya4}
\usepackage{verbatim}

\ifpdf
\usepackage[pdftex]{graphicx}
\pdfcompresslevel=9
\usepackage[bookmarksnumbered=true,pdfpagemode=UseOutlines,
pdfstartview=,bookmarksopen=true]{hyperref}
\else
\usepackage{graphicx}
\usepackage[hypertex,colorlinks]{hyperref}
%\usepackage[dvips]{hyperref}
\fi

\parindent0cm
\parskip1.5ex plus0.5ex minus0.5ex

\begin{document}
\title{EP-Tool library interface}
\author{Markus Multrus, Ralph Sperschneider\\
  Fraunhofer IIS-A}
\maketitle

\section{About the EP-Tool library interface}
The EP-Tool library interface is based on the epTool and provides the
unequal error protection (UEP) capability to the ISO/IEC 14496-3
codecs. The software is provided in source form: 
\begin{description}
\item[libeptool.h] This file contains all the exported routines used by
  the libeptool library. It has to be included in your
  code (\verb?#include "libeptool.h"?).
\item[libeptool.a] This file contains all required functions. You will
  need to link it to your executable.
\end{description}

\section{Namespace/runtime issues}
All exported routines start with the prefix EPTool. The library
does not create any global variables.

\section{Important known limitations of the library}
The library should only be used for the following operations:
\begin{itemize}
\item Converting MPEG-4 audio compressed data provided in
  (epConfig==1) style to data making use of (epConfig==2) style. This
  procedure is referred to as encoding.
\item  Converting MPEG-4 audio compressed data provided in
  (epConfig==2) style to data making use of (epConfig==1) style. This
  procedure is referred to as decoding.
\end{itemize}

MPEG-4 Audio compressed data of all audio object types permitting use of
(epConfig==2) style are supported. The conversion is done on a frame
basis in case of non-scalable object types. The conversion is done on a
layer basis in case of scalable object types. Note:
\begin{itemize}
\item The (epConfig==1) style frames/layers are subdivided in a couple
  of access units, where each access unit contains a single error
  sensitivity category.
\item The (epConfig==2) style frames/layers consist always of one access
  unit.
\end{itemize}

If bitstuffing is disabled, the in- and outputbuffers have to contain
bits, not bytes. Please also enable the bitFlag. 
%The library should only be used with the following limitations: 
%\begin{itemize}
%\item max.~64 classes
%\end{itemize}

\section{Encoding (fundamental usage)}
\begin{description}
\item[Initialization:] 
 EPTool\_CreateExpl()
\item[Frame/Layer encoding:] EPTool\_EncodeFrame()
\item[Termination:] 
  EPTool\_Destroy()
\item[detailed informations:]
    There are two ways, to get the Error Protection Specific
    Configuration (EP Specific Configuration):
    \begin{itemize}
      \item Reading from a file:
        \begin{itemize}
          \item handing over EP Specific Configuration to configBuffer:
          (Syntax of configBuffer as in ISO/IEC 14496-3:1999/Amd.1:2000)

          \begin{itemize}
          \item EPTool\_CreateBufferEPSpecificConfig( unsigned char
          **configBuffer, const char *predFileName );
          \item EPTool\_CreateExpl( HANDLE\_EP\_CONFIG *hEpConfig,  const char
          *predFileName, unsigned char **configBuffer, long
          *configLength, int reorderFlag, int bitFlag );\\
          Ep Specific Configuration now is written to configBuffer, it
          can be handed over to other interfaces.
          \item EPTool\_EncodeFrame (HANDLE\_EP\_CONFIG hEpConfig,
          HANDLE\_CLASS\_FRAME inBuffer, unsigned char *outBuffer, int
          *outLength , int bitFlag );
          \item EPTool\_Destroy( HANDLE\_EP\_CONFIG hEpConfig );
          \item EPTool\_DestroyBufferEPSpecificConfig( unsigned char
            **configBuffer )
          \end{itemize}
          
          \item not handing over EP Specific Configuration in to configBuffer:

          \begin{itemize}
          \item EPTool\_CreateExpl( HANDLE\_EP\_CONFIG *hEpConfig,  const char
          *predFileName, NULL, NULL, int reorderFlag, int bitFlag );
          \item EPTool\_EncodeFrame (HANDLE\_EP\_CONFIG hEpConfig,
          HANDLE\_CLASS\_FRAME inBuffer, unsigned char *outBuffer, int
          *outLength , int bitFlag );
          \item EPTool\_Destroy( HANDLE\_EP\_CONFIG hEpConfig );
          \end{itemize}

        \end{itemize}

      \item Getting EP Specific Configuration from configBuffer:
        (Syntax of configBuffer as in ISO/IEC 14496-3:1999/Amd.1:2000)
        \begin{itemize}
           \item EPTool\_CreateExpl( HANDLE\_EP\_CONFIG *hEpConfig, NULL, unsigned char **configBuffer, long
          *configLength, int reorderFlag, int bitFlag );
            \item EPTool\_EncodeFrame (HANDLE\_EP\_CONFIG hEpConfig,
          HANDLE\_CLASS\_FRAME inBuffer, unsigned char *outBuffer, int
          *outLength , int bitFlag );
            \item EPTool\_Destroy( HANDLE\_EP\_CONFIG hEpConfig );
          \end{itemize}
      \end{itemize}
 
\end{description}

\section{Decoding (fundamental usage)}
\begin{description}
\item[Initialization:] EPTool\_CreateExpl()
\item[Frame/Layer decoding:] EPTool\_DecodeFrame() 
\item[Termination:] EPTool\_Destroy()
\item[detailed informations:]
  See 'encoding' (Substitute EPTool\_EncodeFrame() with EPTool\_DecodeFrame()). 
\end{description}


\section{Error protection functions}
\subsection{Content of libeptool.h}
\verbatiminput{../src/libeptool.h}

\end{document}
